% Terjemahan Bahasa Indonesia dari chapter2.tex baris 35--115.
% Sumber: Methods of Algebra, Volume 2, commit
% 9a5803ff77dd3257484cb177f851a73770a59dd3 (CC BY 4.0).
% stable-unit-id: o014.aljabr2.chapter2.abelian-category-definition

% segment-id: o014.li2.chapter2.abelian-category-definition.h001
\section{Definisi Kategori Abelian}\label{sec:Abel-cat-def}
\hypertarget{o014.aljabr2.chapter2.abelian-category-definition}{ }

% segment-id: o014.li2.chapter2.abelian-category-definition.p001
Pada \S~\ref{sec:Ker-Coker} telah diperkenalkan kategori-$\cate{Ab}$ dan
kategori aditif. Kategori abelian yang dikaji dalam bab ini membentuk kelas
khusus kategori aditif.

% segment-id: o014.li2.chapter2.abelian-category-definition.d001
\begin{definisi}[Kategori abelian]
	\index{kategori abelian (abelian category)}
	Misalkan $\mathcal{A}$ merupakan kategori aditif. Jika setiap morfisme dalam
	$\mathcal{A}$ mempunyai kernel dan kokernel, serta semua morfisme tersebut
	ketat (Definisi~\ref{def:strict-morphism}), maka $\mathcal{A}$ disebut
	\emph{kategori abelian}.

	Jika $\mathcal{A}$ juga merupakan kategori $\Bbbk$-linear
	(Definisi~\ref{def:k-linear-cat}), maka kategori tersebut disebut
	\emph{kategori abelian $\Bbbk$-linear}.
\end{definisi}

% segment-id: o014.li2.chapter2.abelian-category-definition.r001
\begin{catatan}\label{rem:Abel-cat-finite-lim}
	Kategori abelian mempunyai semua $\varprojlim$ hingga dan semua
	$\varinjlim$ hingga. Hal ini karena produk hingga dan koproduk hingga ada
	(keduanya merupakan biproduk), sedangkan penyama dan kopenyama juga ada
	(Catatan~\ref{rem:equalizer-Ker}); data tersebut cukup untuk membangun semua
	$\varprojlim$ hingga dan $\varinjlim$ hingga. Lihat
	\cite[Teorema 2.8.3]{Li1}.
\end{catatan}

% segment-id: o014.li2.chapter2.abelian-category-definition.p002
Menurut Contoh~\ref{eg:strict-morphism}, untuk setiap gelanggang $R$, kategori
$R\dcate{Mod}$ merupakan kategori abelian. Dengan memakai $R^{\text{op}}$
sebagai pengganti $R$, terlihat bahwa kategori modul kanan $\cated{Mod}R$ juga
demikian; dengan mengambil kasus khusus $R = \Z$, terlihat bahwa $\cate{Ab}$
merupakan kategori abelian. Sebaliknya, $\cate{Ban}_{\CC}$ yang ditinjau dalam
contoh yang sama bukan kategori abelian, karena $\cate{Ban}_{\CC}$ mempunyai
morfisme yang tidak ketat.

% segment-id: o014.li2.chapter2.abelian-category-definition.pr001
\begin{proposisi}\label{prop:abelian-cat-duality}
	Jika kategori aditif $\mathcal{A}$ merupakan kategori abelian, maka
	$\mathcal{A}^{\opp}$ juga merupakan kategori abelian.
\end{proposisi}
% segment-id: o014.li2.chapter2.abelian-category-definition.pf001
\begin{bukti}
	Kernel dan kokernel saling dual, sedangkan
	Catatan~\ref{rem:im-coim-duality} menunjukkan bahwa dalam $\mathcal{A}$,
	morfisme $\Coim(f) \to \Image(f)$ tidak lain adalah, dalam
	$\mathcal{A}^{\opp}$, morfisme
	$\Coim(f^{\opp}) \to \Image(f^{\opp})$.
\end{bukti}

% segment-id: o014.li2.chapter2.abelian-category-definition.p003
Salah satu cara abstrak untuk membentuk kategori abelian baru dari kategori
abelian yang telah ada ialah dengan mengambil kategori funktor. Pertama-tama
perhatikan bahwa, jika $\mathcal{A}$ merupakan kategori-$\cate{Ab}$, maka untuk
setiap\footnote{Jika tidak ingin $\mathcal{A}^I$ menjadi kategori besar,
anggaplah $I$ kecil. Hal ini tidak membawa dampak substantif di sini.} kategori
$I$,
	kategori funktor $\mathcal{A}^I$ juga secara alami merupakan
	kategori-$\cate{Ab}$. Untuk funktor
	$F, G: I \rightrightarrows \mathcal{A}$, penjumlahan dalam
	$\Hom_{\mathcal{A}^I}(F, G)$ didefinisikan pada setiap objek melalui
$(\varphi_i)_i + (\psi_i)_i := (\varphi_i + \psi_i)_i$.

% segment-id: o014.li2.chapter2.abelian-category-definition.pr002
\begin{proposisi}\label{prop:functor-cat-Abel}
	Misalkan $\mathcal{A}$ merupakan kategori abelian. Untuk setiap kategori
	$I$, kategori funktor $\mathcal{A}^I$ juga merupakan kategori abelian.
\end{proposisi}
% segment-id: o014.li2.chapter2.abelian-category-definition.pf002
\begin{bukti}
	Berdasarkan Proposisi~\ref{prop:forget-create}, produk dan koproduk hingga,
	serta kernel dan kokernel, dalam $\mathcal{A}^I$ dibangun pada setiap objek.
	Citra dan kocitra suatu morfisme, dan karena itu juga morfisme kanonik
	$\Coim \to \Image$, dibangun dengan cara yang sama. Dengan demikian,
	semua syarat kategori abelian bagi $\mathcal{A}^I$ dapat direduksi menjadi
	syarat-syarat yang sama bagi $\mathcal{A}$.
\end{bukti}

% segment-id: o014.li2.chapter2.abelian-category-definition.p004
Setiap morfisme $f: X \to Y$ dalam kategori abelian mempunyai faktorisasi
epi--mono yang unik (Proposisi~\ref{prop:epi-mono-uniqueness}), yang secara
eksplisit dapat diambil sebagai
$X \twoheadrightarrow \left(\Coim(f) \simeq \Image(f)\right)
\hookrightarrow Y$. Karena itu, Lema~\ref{prop:mono-epi-ker-coker}~(iii)
langsung memberikan
% segment-id: o014.li2.chapter2.abelian-category-definition.q001
\[ \Ker(f) = \Ker[X \to \Image(f)], \quad \Coker(f) = \Coker[\Image(f) \to Y]. \]

% segment-id: o014.li2.chapter2.abelian-category-definition.r002
\begin{catatan}\label{rem:Abel-cat-fibered-product}
	Untuk morfisme-morfisme $X \xrightarrow{f} Z \xleftarrow{g} Y$ dalam
	kategori abelian, konstruksi produk serat
	(lihat Catatan~\ref{rem:Abel-cat-finite-lim}) menunjukkan bahwa
	$X \dtimes{Z} Y$ tidak lain adalah penyama dari
	% segment-id: o014.li2.chapter2.abelian-category-definition.g001
	$\begin{tikzcd}
		X \oplus Y \arrow[r, yshift=0.3em, "{(f,0)}"] \arrow[r, yshift=-0.3em, "{(0,g)}"'] & Z
	\end{tikzcd}$,
	yakni $\Ker[ X \oplus Y \xrightarrow{(f, -g)} Z ]$. Tentu saja, morfisme
	$X \leftarrow X \dtimes{Z} Y \rightarrow Y$ yang dibawa oleh produk serat
	berasal dari proyeksi
	$X \xleftarrow{p_1} X \oplus Y \xrightarrow{p_2} Y$.

	Untuk $X' \xleftarrow{f'} Z' \xrightarrow{g'} Y'$, konstruksi koproduk
	serat merealisasikan $X' \dsqcup{Z'} Y'$ sebagai
	$\Coker[ Z' \xrightarrow{(-f', g')} X' \oplus Y']$, sedangkan morfisme
	$X' \rightarrow X' \dsqcup{Z'} Y' \leftarrow Y'$ berasal dari
	$X' \xrightarrow{\iota_1} X' \oplus Y'
	\xleftarrow{\iota_2} Y'$. Sifat-sifat ini memang jelas dalam kasus modul.
\end{catatan}

% segment-id: o014.li2.chapter2.abelian-category-definition.p005
Hasil berikut akan digunakan berulang kali.

% segment-id: o014.li2.chapter2.abelian-category-definition.pr003
\begin{proposisi}\label{prop:Abel-cat-pull-push}
	Misalkan diagram dalam kategori abelian $\mathcal{A}$
	% segment-id: o014.li2.chapter2.abelian-category-definition.g002
	\[\begin{tikzcd}
		X \arrow[r, "f"] \arrow[d, "a"'] & Y \arrow[d, "b"] \\
		X' \arrow[r, "g"'] & Y'
	\end{tikzcd}\]
	merupakan diagram tarik balik (atau dorong keluar), dan $g$ merupakan
	epimorfisme (atau $f$ merupakan monomorfisme). Maka diagram tersebut juga
	merupakan diagram dorong keluar (atau tarik balik), dan $f$ merupakan
	epimorfisme (atau $g$ merupakan monomorfisme).
\end{proposisi}
% segment-id: o014.li2.chapter2.abelian-category-definition.pf003
\begin{bukti}
	Cukup ditangani kasus diagram tarik balik. Akan ditunjukkan bahwa
	% segment-id: o014.li2.chapter2.abelian-category-definition.l001
	\begin{compactitem}
		% segment-id: o014.li2.chapter2.abelian-category-definition.i001
		\item $X \xrightarrow{(a, f)} X' \oplus Y$ memberikan
		$\Ker\left[ X' \oplus Y \xrightarrow{(g, -b)} Y' \right]$;
		% segment-id: o014.li2.chapter2.abelian-category-definition.i002
		\item $(g, -b): X' \oplus Y \to Y'$ merupakan epimorfisme.
	\end{compactitem}

	Pernyataan mengenai $\Ker$ langsung mengikuti deskripsi
	$X \simeq X' \dtimes{Y'} Y$ sebagai kernel dalam
	Catatan~\ref{rem:Abel-cat-fibered-product}. Selain itu, karena
	$g = (g, -b) \iota_1$, sifat epimorfisme $g$ mengakibatkan
	$(g, -b)$ juga merupakan epimorfisme.

	Dengan demikian,
	$Y' = \Image((g, -b)) \leftiso
	\Coker[X \xrightarrow{(a,f)} X' \oplus Y ]$. Dengan mengubah sudut pandang
	dan memakai deskripsi koproduk serat sebagai kokernel dalam
	Catatan~\ref{rem:Abel-cat-fibered-product}, segera terlihat bahwa
	% segment-id: o014.li2.chapter2.abelian-category-definition.g003
	\[\begin{tikzcd}
		X \arrow[r, "f"] \arrow[d, "-a"'] & Y \arrow[d, "-b"] \\
		X' \arrow[r, "g"'] & Y'
	\end{tikzcd}\]
	merupakan diagram dorong keluar. Dengan mengomposisikan masing-masing kolom
	dengan automorfisme $-\identity_X$ dan $-\identity_Y$,
	diperoleh kembali diagram semula, yang tetap merupakan diagram dorong
	keluar. Dalam kasus dorong keluar, telah diketahui bahwa
	$\Coker(f) \rightiso \Coker(g)$
	(Proposisi~\ref{prop:pull-back-ker}); karena itu, $f$ merupakan epimorfisme
	(Lema~\ref{prop:mono-epi-ker-coker}).
\end{bukti}

% segment-id: o014.li2.chapter2.abelian-category-definition.co001
\begin{korolari}\label{prop:Abel-cat-subquotient}
	Subkuosien dalam kategori abelian dapat didefinisikan secara ekuivalen
	sebagai objek kuosien dari suatu subobjek, atau sebagai subobjek dari suatu
	objek kuosien.
\end{korolari}
% segment-id: o014.li2.chapter2.abelian-category-definition.pf004
\begin{bukti}
	Proposisi~\ref{prop:Abel-cat-pull-push} menunjukkan bahwa
	Proposisi~\ref{prop:subquotient-equiv} dapat diterapkan pada setiap kategori
	abelian.
\end{bukti}

% segment-id: o014.li2.chapter2.abelian-category-definition.p006
Untuk sebarang gelanggang $R$, semua hasil di atas sudah lazim bagi kategori
modul kiri $R\dcate{Mod}$; hal yang sama berlaku bagi kategori modul kanan.
